\section{Research Gap and Context}
\label{sec:research_gap}

Despite their promise, production ZK-rollups (such as zkSync Era, Starknet, and Polygon zkEVM) are highly complex systems. Their performance is dictated by an intricate interplay of several bottlenecks: sequencer scheduling policies, batch formation triggers, data availability (DA) modes, L1 submission costs, and highly compute-bound proof generation overhead. The complexity of production systems often makes it difficult to conduct controlled, isolated experiments to answer fundamental questions about these interactions. 

Existing benchmarking efforts typically rely on either purely theoretical models or observational analysis of active production networks. Because production L2 networks are optimized for live deployment, reliability, and user security, they actively resist controlled benchmarking. Forking a production rollup to study parameters introduces too much noise from custom compilers and complex gas logic, making layer-level variables impossible to isolate. Furthermore, Ethereum's use of Hexary Patricia Tries (PMT) for state management is notoriously slow to prove in ZK due to variable path lengths and complex serialization. 

This report details the development and evaluation of a clean-room modular ZK-Rollup benchmarking prototype, developed as a final year research project. The prototype is explicitly a \textit{research prototype} and is not intended to be a production-ready mainnet rollup or a complete zkEVM implementation. By implementing a 256-level binary Sparse Merkle Tree (SMT) instead of a PMT, the prototype guarantees simple, constant-length ZK proofs. Instead of executing full EVM bytecode, it serves as a controlled benchmarking framework designed specifically to isolate and evaluate the impact of different configuration choices on system performance using synthetic transfer-centric workloads.

By cleanly separating sequencing, execution, proof handling, submission, and data availability concerns, the architecture enables reproducible experiments. We investigate how varying batch sizes, scheduling policies, and data availability mechanisms (such as standard calldata versus EIP-4844 blobs) influence the critical throughput--latency--cost trade-off. Furthermore, the system incorporates a benchmark-mode mock prover that introduces parameterized proof delays, allowing us to analytically model proof overhead without conflating systems-level queueing dynamics with the massive hardware requirements of actual cryptographic proof generation, while still optionally supporting real RISC Zero STARK proofs to measure cycle limits.
